\section{Limitations, Implications, and Conclusion}
\label{sec:limitations}

\paragraph{What is established.}
Candidate-indexed mass can make a shared-label active experiment depend on roster representation. Exact-quality T2-A and fixed-query equality establish the complete refinement-to-root-cost causal chain, and learned raw-input adapters show that holdout-reference lookup is not necessary for a terminal effect. The theory explains why responses alone cannot generally recover the missing frame.

\paragraph{What is not established.}
The natural anchor is one pair and causes no reliable terminal harm. T2-A uses public or independently acquired references; T1 passes $3/8$, LLM Selector $2/6$, and untargeted additions are mostly benign. Most learned same-$s$ tasks are terminal-null or stop. IMDB is disjoint held-out evidence, but its condition was developed on 12,000 rows of the same benchmark; prospective Yelp misses its $.5$-point materiality gate, and both are binary sentiment. The adapters learn separate weights while sharing public base architectures, and no endpoint is admitted to a live registry. We therefore do not establish production compromise, clinical harm, universal terminal damage, adversarial intent, independent full-base checkpoints, or the first sequential clone effect.

\paragraph{Implication.}
The operational requirement is not a universal deduplication metric. An evaluator must declare what a root means for its claim, bind entries to that unit using appropriate provenance or governance evidence, assign fixed root mass, and specify a refinement-invariant within-root summary (Proposition~\ref{prop:factorization}). When authentication is unavailable, multiplicity sensitivity is part of the uncertainty of the reported selection result and should be audited explicitly.
